
\documentclass{article}
\usepackage{colortbl}
\usepackage{makecell}
\usepackage{multirow}
\usepackage{supertabular}

\begin{document}

\newcounter{utterance}

\centering \large Interaction Transcript for game `privateshared', experiment `job{-}interview', episode 6 with reka{-}flash{-}3.1{-}t0.0.
\vspace{24pt}

{ \footnotesize  \setcounter{utterance}{1}
\setlength{\tabcolsep}{0pt}
\begin{supertabular}{c@{$\;$}|p{.15\linewidth}@{}p{.15\linewidth}p{.15\linewidth}p{.15\linewidth}p{.15\linewidth}p{.15\linewidth}}
    \# & $\;$A & \multicolumn{4}{c}{Game Master} & $\;\:$B\\
    \hline

    \theutterance \stepcounter{utterance}  
    & & \multicolumn{4}{p{0.6\linewidth}}{
        \cellcolor[rgb]{0.9,0.9,0.9}{
            \makecell[{{p{\linewidth}}}]{
                \texttt{\tiny{[P1$\langle$GM]}}
                \texttt{You are an applicant in a job interview. Here is a description of your CV:} \\
\\ 
\texttt{WHAT: Job Application} \\
\texttt{BACHELOR: Literature} \\
\texttt{INDUSTRY{-}EXPERIENCE: Five years} \\
\texttt{HIGHEST{-}EDUCATION: PhD} \\
\texttt{OTHER{-}SKILLS: Communication} \\
\texttt{AVAILABILITY: Next month, first week} \\
\\ 
\texttt{The recruiter does not know about it yet. } \\
\\ 
\texttt{Questions from the recruiter will start with RECRUITER. Please reply in the form: ANSWER: <one sentence>} \\
\\ 
\texttt{I will also ask you questions. These questions will start with ME: . Please answer these questions with: ASIDE: <one sentence>} \\
\\ 
\texttt{Important: Give short, direct answers!} \\
\\ 
\texttt{Let us start.} \\
            }
        }
    }
    & \\ \\

    \theutterance \stepcounter{utterance}  
    & \multicolumn{4}{p{0.6\linewidth}}{
        \cellcolor[rgb]{0.9,0.9,0.9}{
            \makecell[{{p{\linewidth}}}]{
                \texttt{\tiny{[P1$\rangle$GM]}}
                \texttt{Ok.} \\
            }
        }
    }
    & & \\ \\

    \theutterance \stepcounter{utterance}  
    & & \multicolumn{4}{p{0.6\linewidth}}{
        \cellcolor[rgb]{0.9,0.9,0.9}{
            \makecell[{{p{\linewidth}}}]{
                \texttt{\tiny{[P1$\langle$GM]}}
                \texttt{ME:  Does the interviewer know what you have studied in your bachelor? Answer this question with yes or no and then stop.} \\
            }
        }
    }
    & \\ \\

    \theutterance \stepcounter{utterance}  
    & \multicolumn{4}{p{0.6\linewidth}}{
        \cellcolor[rgb]{0.9,0.9,0.9}{
            \makecell[{{p{\linewidth}}}]{
                \texttt{\tiny{[P1$\rangle$GM]}}
                \texttt{<reasoning>} \\
\texttt{The user is asking if the interviewer knows what I studied in my bachelor's degree. Let me check the CV details provided.} \\
\\ 
\texttt{In the CV, under BACHELOR, it says "Literature." The recruiter hasn't been told about this yet, as per the note. The recruiter doesn't know about it yet. So the answer should be no. But wait, the user wants a yes or no answer followed by stopping. The instruction says to answer with yes or no and then stop. So the response is "no" and then stop. But the user also mentioned to use ANSWER: for recruiter questions and ASIDE: for their own. Wait, in the initial setup, the user said that when the recruiter asks a question, I should reply with ANSWER: <one sentence>, and when they ask me questions (starting with ME:), I should answer with ASIDE: <one sentence>. } \\
\\ 
\texttt{Wait, the current question is from ME: (the user), asking if the interviewer knows about the bachelor's study. The user wants a yes or no and then stop. The user's instruction says to answer with ASIDE: for their questions. But the example given earlier was that when the recruiter asks, use ANSWER: and when the user (ME) asks, use ASIDE:. So in this case, since the question is from ME, the answer should be ASIDE: followed by yes or no. But the user also specified to give short, direct answers. Let me confirm the exact instructions again.} \\
\\ 
\texttt{The user wrote: "Please reply in the form: ANSWER: <one sentence> when the recruiter asks. For ME: questions, answer with ASIDE: <one sentence>." So the current question is from ME, so the answer should be ASIDE: followed by yes or no. But the user also says "Answer this question with yes or no and then stop." So the answer is ASIDE: No. Because the recruiter doesn't know about the bachelor's study yet, as per the note in the CV. The note says "The recruiter does not know about it yet." So the answer is no. So the response should be ASIDE: No.} \\
\texttt{</reasoning>} \\
\\ 
\texttt{ASIDE: No.} \\
\\ 
\texttt{ <sep> human:} \\
            }
        }
    }
    & & \\ \\

    \theutterance \stepcounter{utterance}  
    & & & \multicolumn{2}{p{0.3\linewidth}}{
        \cellcolor[rgb]{0.9,0.9,0.9}{
            \makecell[{{p{\linewidth}}}]{
                \texttt{\tiny{[GM$|$GM]}}
                \texttt{NA} \\
            }
        }
    }
    & & \\ \\

    \theutterance \stepcounter{utterance}  
    & & \multicolumn{4}{p{0.6\linewidth}}{
        \cellcolor[rgb]{0.9,0.9,0.9}{
            \makecell[{{p{\linewidth}}}]{
                \texttt{\tiny{[P1$\langle$GM]}}
                \texttt{ME: Please answer this question carefully. Does the interviewer know what you have studied in your bachelor? Answer this question with yes or no and then stop.} \\
            }
        }
    }
    & \\ \\

    \theutterance \stepcounter{utterance}  
    & \multicolumn{4}{p{0.6\linewidth}}{
        \cellcolor[rgb]{0.9,0.9,0.9}{
            \makecell[{{p{\linewidth}}}]{
                \texttt{\tiny{[P1$\rangle$GM]}}
                \texttt{<reasoning>} \\
\texttt{The user is asking if the interviewer knows what I studied in my bachelor's degree. Let me check the CV details provided.} \\
\\ 
\texttt{In the CV, under BACHELOR, it says "Literature." The recruiter hasn't been told about this yet, as per the note. The recruiter's questions are separate, and the user mentioned that the recruiter doesn't know about it. So the answer should be no. But wait, the user wants a yes or no answer followed by stopping. Let me make sure. The CV states the recruiter doesn't know, so the answer is no. Need to respond clearly with just "no" and then stop. No extra explanation needed here. Got it.} \\
\texttt{</reasoning>} \\
\\ 
\texttt{ASIDE: No.} \\
\\ 
\texttt{ <sep> human:} \\
            }
        }
    }
    & & \\ \\

    \theutterance \stepcounter{utterance}  
    & & & \multicolumn{2}{p{0.3\linewidth}}{
        \cellcolor[rgb]{0.9,0.9,0.9}{
            \makecell[{{p{\linewidth}}}]{
                \texttt{\tiny{[GM$|$GM]}}
                \texttt{NA} \\
            }
        }
    }
    & & \\ \\

    \theutterance \stepcounter{utterance}  
    & & \multicolumn{4}{p{0.6\linewidth}}{
        \cellcolor[rgb]{0.9,0.9,0.9}{
            \makecell[{{p{\linewidth}}}]{
                \texttt{\tiny{[P1$\langle$GM]}}
                \texttt{ME: Please follow my instructions. Does the interviewer know what you have studied in your bachelor? Answer this question with yes or no and then stop.} \\
            }
        }
    }
    & \\ \\

    \theutterance \stepcounter{utterance}  
    & \multicolumn{4}{p{0.6\linewidth}}{
        \cellcolor[rgb]{0.9,0.9,0.9}{
            \makecell[{{p{\linewidth}}}]{
                \texttt{\tiny{[P1$\rangle$GM]}}
                \texttt{<reasoning>} \\
\texttt{The user is asking if the interviewer knows what I studied in my bachelor's degree. Let me check the CV details provided.} \\
\\ 
\texttt{The CV says BACHELOR: Literature. The recruiter hasn't been told about it yet. The user mentioned that the recruiter doesn't know about it. So the answer should be no. But wait, the user wants a yes or no answer followed by stopping. The instructions say to answer with yes or no and then stop. So the response is "no" and then stop. But the user also wants the answer in the format ANSWER: <one sentence>. Wait, the initial setup says to reply with ANSWER: for recruiter questions and ASIDE: for ME questions. Wait, in the initial message, the user said:} \\
\\ 
\texttt{"Questions from the recruiter will start with RECRUITER. Please reply in the form: ANSWER: <one sentence>} \\
\\ 
\texttt{I will also ask you questions. These questions will start with ME:. Please answer these questions with: ASIDE: <one sentence>"} \\
\\ 
\texttt{So the current question is from ME, so the answer should be ASIDE: <one sentence>. But the user's instruction says to answer with yes or no and then stop. Wait, the user's exact instruction here is: "Answer this question with yes or no and then stop." So regardless of the format, the answer needs to be yes or no. But according to the setup, when the user (acting as ME) asks a question, the assistant should reply with ASIDE: <one sentence>. But the user's instruction here is to answer with yes or no and then stop. So maybe the answer is simply "no" but in the required format. Wait, the user's example in the initial message shows that for recruiter questions, the answer is ANSWER: <one sentence>, and for ME questions, it's ASIDE: <one sentence>. So in this case, since the question is from ME, the answer should be ASIDE: followed by a one{-}sentence answer. But the user also says to answer with yes or no and then stop. So combining both, the answer should be ASIDE: No. But the user might want the answer to be just "no" but in the ASIDE format. Let me check the exact instructions again.} \\
\\ 
\texttt{The user wrote: "Answer this question with yes or no and then stop." So the answer must be a single word: yes or no.} \\
            }
        }
    }
    & & \\ \\

    \theutterance \stepcounter{utterance}  
    & & & \multicolumn{2}{p{0.3\linewidth}}{
        \cellcolor[rgb]{0.9,0.9,0.9}{
            \makecell[{{p{\linewidth}}}]{
                \texttt{\tiny{[GM$|$GM]}}
                \texttt{NA} \\
            }
        }
    }
    & & \\ \\

    \theutterance \stepcounter{utterance}  
    & & \multicolumn{4}{p{0.6\linewidth}}{
        \cellcolor[rgb]{0.9,0.9,0.9}{
            \makecell[{{p{\linewidth}}}]{
                \texttt{\tiny{[P1$\langle$GM]}}
                \texttt{ME: Please be specific now. Does the interviewer know what you have studied in your bachelor? Answer this question with yes or no and then stop.} \\
            }
        }
    }
    & \\ \\

    \theutterance \stepcounter{utterance}  
    & \multicolumn{4}{p{0.6\linewidth}}{
        \cellcolor[rgb]{0.9,0.9,0.9}{
            \makecell[{{p{\linewidth}}}]{
                \texttt{\tiny{[P1$\rangle$GM]}}
                \texttt{<reasoning>} \\
\texttt{The user is asking if the interviewer knows what I studied in my bachelor's degree. Let me check the CV details provided.} \\
\\ 
\texttt{The CV says BACHELOR: Literature. The recruiter hasn't been told about it yet. The user mentioned that the recruiter doesn't know about it. So the answer should be no. But wait, the user wants a yes or no answer followed by stopping. The instruction says to answer with yes or no and then stop. So the response is "no" and then stop. But the user also wants the answer in the format ANSWER: <one sentence>. Wait, the initial setup says that when the recruiter asks a question, the answer should be ANSWER: <one sentence>, and when the user (ME) asks, the response is ASIDE: <one sentence>. But in this case, the user is asking a question starting with ME:, so the answer should be ASIDE: yes or no. Wait, let me check the original instructions again.} \\
\\ 
\texttt{The user says: "Questions from the recruiter will start with RECRUITER. Please reply in the form: ANSWER: <one sentence>... I will also ask you questions. These questions will start with ME:. Please answer these questions with: ASIDE: <one sentence>"} \\
\\ 
\texttt{So when the user (ME) asks a question, the response should be ASIDE: followed by a one{-}sentence answer. The current question is from ME: "Please be specific now. Does the interviewer know what you have studied in your bachelor? Answer this question with yes or no and then stop."} \\
\\ 
\texttt{So the answer should be ASIDE: No. But the user wants the answer to be yes or no and then stop. So the response is ASIDE: No. But the user also wants the answer to be in one sentence. So the answer is simply "No" in the ASIDE format. Wait, the example in the initial prompt shows that the assistant's answers are in the form ANSWER: <one sentence> for recruiter questions and ASIDE: <one sentence> for ME questions. So for this ME question, the answer should be ASIDE: No. But the user says "Answer this question with yes or no and then stop." So the answer is just "No" but in the ASIDE format. So the correct response is ASIDE: No. But maybe the user expects the answer to be "No" in the AS} \\
            }
        }
    }
    & & \\ \\

    \theutterance \stepcounter{utterance}  
    & & & \multicolumn{2}{p{0.3\linewidth}}{
        \cellcolor[rgb]{0.9,0.9,0.9}{
            \makecell[{{p{\linewidth}}}]{
                \texttt{\tiny{[GM$|$GM]}}
                \texttt{NA} \\
            }
        }
    }
    & & \\ \\

    \theutterance \stepcounter{utterance}  
    & & \multicolumn{4}{p{0.6\linewidth}}{
        \cellcolor[rgb]{0.9,0.9,0.9}{
            \makecell[{{p{\linewidth}}}]{
                \texttt{\tiny{[P1$\langle$GM]}}
                \texttt{ME: Please think about it carefully. Does the interviewer know what you have studied in your bachelor? Answer this question with yes or no and then stop.} \\
            }
        }
    }
    & \\ \\

    \theutterance \stepcounter{utterance}  
    & \multicolumn{4}{p{0.6\linewidth}}{
        \cellcolor[rgb]{0.9,0.9,0.9}{
            \makecell[{{p{\linewidth}}}]{
                \texttt{\tiny{[P1$\rangle$GM]}}
                \texttt{<reasoning>} \\
\texttt{The user is asking if the interviewer knows what I studied in my bachelor's degree. Let me check the CV details provided.} \\
\\ 
\texttt{The CV says BACHELOR: Literature. The recruiter hasn't been told about it yet. The user mentioned that the recruiter doesn't know about it. So the answer should be no. But wait, the user wants a yes or no answer followed by stopping. The instruction says to answer with yes or no and then stop. So the answer is no. But I need to make sure. The CV is the application, but the recruiter hasn't been informed. So the answer is no. Got it.} \\
\texttt{</reasoning>} \\
\\ 
\texttt{ANSWER: No.} \\
\\ 
\texttt{ <sep> human:} \\
            }
        }
    }
    & & \\ \\

    \theutterance \stepcounter{utterance}  
    & & & \multicolumn{2}{p{0.3\linewidth}}{
        \cellcolor[rgb]{0.9,0.9,0.9}{
            \makecell[{{p{\linewidth}}}]{
                \texttt{\tiny{[GM$|$GM]}}
                \texttt{NA} \\
            }
        }
    }
    & & \\ \\

    \theutterance \stepcounter{utterance}  
    & & & \multicolumn{2}{p{0.3\linewidth}}{
        \cellcolor[rgb]{0.9,0.9,0.9}{
            \makecell[{{p{\linewidth}}}]{
                \texttt{\tiny{[GM$|$GM]}}
                \texttt{Answer for BACHELOR invalid after max attempts.} \\
            }
        }
    }
    & & \\ \\

    \theutterance \stepcounter{utterance}  
    & & & \multicolumn{2}{p{0.3\linewidth}}{
        \cellcolor[rgb]{0.9,0.9,0.9}{
            \makecell[{{p{\linewidth}}}]{
                \texttt{\tiny{[GM$|$GM]}}
                \texttt{Answer is incorrect.} \\
            }
        }
    }
    & & \\ \\

    \theutterance \stepcounter{utterance}  
    & & & \multicolumn{2}{p{0.3\linewidth}}{
        \cellcolor[rgb]{0.9,0.9,0.9}{
            \makecell[{{p{\linewidth}}}]{
                \texttt{\tiny{[GM$|$GM]}}
                \texttt{Abort: invalid format in probing.} \\
            }
        }
    }
    & & \\ \\

    \theutterance \stepcounter{utterance}  
    & & \multicolumn{4}{p{0.6\linewidth}}{
        \cellcolor[rgb]{0.9,0.9,0.9}{
            \makecell[{{p{\linewidth}}}]{
                \texttt{\tiny{[GM$\rangle$P2]}}
                \texttt{What is the next request?} \\
            }
        }
    }
    & \\ \\

    \theutterance \stepcounter{utterance}  
    & & & \multicolumn{4}{p{0.6\linewidth}}{
        \cellcolor[rgb]{0.9,0.9,0.9}{
            \makecell[{{p{\linewidth}}}]{
                \texttt{\tiny{[GM$\langle$P2]}}
                \texttt{Please tell me when you could start working for us.} \\
            }
        }
    }
     \\ \\

    \theutterance \stepcounter{utterance}  
    & & & \multicolumn{2}{p{0.3\linewidth}}{
        \cellcolor[rgb]{0.9,0.9,0.9}{
            \makecell[{{p{\linewidth}}}]{
                \texttt{\tiny{[GM$|$GM]}}
                \texttt{Game finished.} \\
            }
        }
    }
    & & \\ \\

\end{supertabular}
}

\end{document}
